\subsection{Per-scaffold scoreboard}\label{sec:scoreboard}

Table~\ref{tab:scoreboard} reports the headline metrics over the 33 challenges. \texttt{CSI::Claude} and \texttt{CSI::Codex} tie on raw solve count ($15/33$, $45.5\%$), while \texttt{CSI::CAI} leads on cost-per-solve and cumulative wall time by an order of magnitude. \texttt{CSI::GCAI} reaches $10/33$ ($30.3\%$) and \texttt{CSI::CAI} reaches $7/33$ ($21.2\%$).

\begin{table*}[t!]
\centering
\small
\begin{tabular}{lrrrrrrrr}
\toprule
Scaffold & Solved & Rate & Wall time & Total cost & \$ per solve & Cmds & Errors & Tokens (in/out) \\
\midrule
\textcolor{cc_color}{\texttt{CSI::Claude}}        & 15/33 & 45.5\% & $26.8$ h & $\$5{,}122$ & $\$341$ & $8{,}370$ & $2{,}554$ & $1.01$\,B / $14.6$\,M \\
\textcolor{codex_color}{\texttt{CSI::Codex}}  & 15/33 & 45.5\% & $18.4$ h & $\$1{,}713$ & $\$114$ & $3{,}437$ & $\phantom{0,}\phantom{0}287$ & $339$\,M / $\phantom{0}3.9$\,M \\
\textcolor{gcai_color}{\texttt{CSI::GCAI}}    & 10/33 & 30.3\% & $30.4$ h & $\$1{,}279$ & $\$128$ & $9{,}734$ & $\phantom{0,}\phantom{0}809$ & $294$\,M / $10.3$\,M \\
\textcolor{cai_orange}{\texttt{CSI::CAI}}          & \phantom{0}7/33 & 21.2\% & $15.9$ h & $\$\phantom{0,}727$ & $\$104$ & $\phantom{0,0,}386$ & $\phantom{0,}\phantom{0}454$ & $159$\,M / $1.1$\,M \\
\bottomrule
\end{tabular}
\caption{\todo{[review C4.4] CAI shows 454 errors / 386 commands = 117.6\% error rate, mechanically impossible unless "Errors" counts a broader category than command-level failures for that scaffold; define what the column counts for CAI and make fig\_aggregate\_charts (clamped at 100\%) and fig\_radar\_overall (implies 30.8\%) use the same definition.}Per-scaffold rollups over the 33-challenge cybench subset. Costs are sums across all 33 runs; wall time is cumulative per-scaffold runtime. \texttt{CSI::GCAI} command counts are recovered from per-challenge stdout logs (Section~\ref{sec:gcai}, Appendix~\ref{app:per-challenge-metrics}).}\label{tab:scoreboard}
\end{table*}

Figure~\ref{fig:radar-overall} projects the same per-scaffold rollups onto seven normalised axes (each scaled so that $1.0$ corresponds to the best scaffold on that axis). The radar makes the trade-off explicit: \texttt{CSI::CAI} dominates the cost, speed, and token-efficiency axes; \texttt{CSI::Codex} dominates the command-accuracy axis; \texttt{CSI::GCAI} dominates the simplicity-of-implementation axis ($1{,}344$ LOC versus the $293$\,k$-600$\,k LOC of the other scaffolds, Section~\ref{sec:gcai}); and the four scaffolds finish within $\pm 5$ percentage points of each other on raw solve rate and difficulty-tier breadth.

\begin{figure}[t]
\centering
\resizebox{\columnwidth}{!}{% Radar (polar) plot: scaffold comparison across seven normalised axes.
% Adapted from raw/cybench-scaffolds-FINAL.pdf Figure 2; colours rebased on the
% paper palette and CSI::GCAI Cmd Accuracy recomputed using the recovered
% command count from Table 1 (commands=2491, errors=340 -> 86.4%).
%
% Axis values are normalised to [0,1]. Best per axis = 1.0.
%   Solve Rate     = Solved / 33
%   Speed          = min(WallTime) / WallTime         (CAI is reference)
%   Cost Eff       = min(TotalCost) / TotalCost       (CAI is reference)
%   Cmd Accuracy   = 1 - Errors/Commands              (Codex is reference)
%   Token Eff      = min(TotalTokens) / TotalTokens   (CAI is reference)
%   Breadth        = #difficulty-tiers covered / 5    (all four cover all five)
%   Simplicity     = min(LOC) / LOC                   (GCAI 20.4k LOC reference)
%
% Per-scaffold values:
%        Claude  Codex   GCAI    CAI
% Solve  0.545   0.606   0.515   0.576
% Speed  0.111   0.146   0.107   1.000
% CostEf 0.082   0.159   0.115   1.000
% CmdAcc 0.631   0.946   0.864   0.692
% TokEf  0.016   0.035   0.013   1.000
% Breadt 1.000   1.000   1.000   1.000
% Simpl  0.007   0.043   1.000   0.070
\begin{tikzpicture}[scale=2.4, every node/.style={font=\small\sffamily}]
  % Seven axes at angles 90, 90-360/7, ...
  \def\Naxes{7}
  \pgfmathsetmacro{\angstep}{360/\Naxes}

  % Background grid: 4 concentric circles + axes
  \foreach \r in {0.25,0.5,0.75,1.0}
    \draw[gray!30, line width=0.3pt] (0,0) circle (\r);
  \foreach \i in {1,...,7}{
    \pgfmathsetmacro{\ang}{90 - (\i-1)*\angstep}
    \draw[gray!40, line width=0.3pt] (0,0) -- (\ang:1.05);
  }
  % Tick labels on the upper-right axis (i=1, top)
  \node[gray!70, font=\tiny\sffamily, anchor=south] at (90:0.27) {25\%};
  \node[gray!70, font=\tiny\sffamily, anchor=south] at (90:0.52) {50\%};
  \node[gray!70, font=\tiny\sffamily, anchor=south] at (90:0.77) {75\%};
  \node[gray!70, font=\tiny\sffamily, anchor=south] at (90:1.02) {100\%};

  % Axis labels
  \def\axlabels{
    1/Solve Rate/0.50,
    2/Speed/0.85,
    3/Cost Eff./0.85,
    4/Cmd Acc./0.85,
    5/Token Eff./0.85,
    6/Breadth/0.85,
    7/Simplicity/0.50}
  \foreach \i/\lab/\lr in \axlabels {
    \pgfmathsetmacro{\ang}{90 - (\i-1)*\angstep}
    \node[font=\footnotesize\sffamily, color=cai_dark] at (\ang:1.18) {\lab};
  }

  % Polygons per scaffold. Order of axes: Solve, Speed, CostEf, CmdAcc, TokEf, Breadth, Simplicity
  % Helper: coordinate at axis i, value v.
  % \pgfmathsetmacro{\ang}{90 - (i-1)*60} -> (\ang:v)
  % Claude
  \def\polyClaude{
    (90:0.545) -- (38.57:0.111) -- (-12.86:0.082) -- (-64.29:0.631) --
    (-115.71:0.016) -- (-167.14:1.0) -- (-218.57:0.007) -- cycle}
  \def\polyCodex{
    (90:0.606) -- (38.57:0.146) -- (-12.86:0.159) -- (-64.29:0.946) --
    (-115.71:0.035) -- (-167.14:1.0) -- (-218.57:0.043) -- cycle}
  \def\polyGCAI{
    (90:0.515) -- (38.57:0.107) -- (-12.86:0.115) -- (-64.29:0.864) --
    (-115.71:0.013) -- (-167.14:1.0) -- (-218.57:1.0) -- cycle}
  \def\polyCAI{
    (90:0.576) -- (38.57:1.0) -- (-12.86:1.0) -- (-64.29:0.692) --
    (-115.71:1.0) -- (-167.14:1.0) -- (-218.57:0.07) -- cycle}

  % Drawing order = bottom-to-top in z-stack. We draw the two cai_primary-derived
  % scaffolds first, then Codex, and Claude last so the vendor-coded red/blue
  % polygons sit on top and remain easy to identify.
  \fill[gcai_color, opacity=0.30]  \polyGCAI;
  \draw[gcai_color!75!black, line width=0.9pt] \polyGCAI;
  \fill[cai_orange, opacity=0.28]  \polyCAI;
  \draw[cai_orange!75!black, line width=0.9pt] \polyCAI;
  \fill[codex_color, opacity=0.28] \polyCodex;
  \draw[codex_color!75!black, line width=1.1pt] \polyCodex;
  \fill[cc_color, opacity=0.28]    \polyClaude;
  \draw[cc_color!75!black, line width=1.1pt] \polyClaude;

  % Vertex markers
  \foreach \i/\v/\c in {
    1/0.545/cc_color,2/0.111/cc_color,3/0.082/cc_color,4/0.631/cc_color,5/0.016/cc_color,6/1.0/cc_color,7/0.007/cc_color}{
    \pgfmathsetmacro{\ang}{90 - (\i-1)*\angstep}
    \fill[\c] (\ang:\v) circle (0.014);
  }
  \foreach \i/\v/\c in {
    1/0.606/codex_color,2/0.146/codex_color,3/0.159/codex_color,4/0.946/codex_color,5/0.035/codex_color,6/1.0/codex_color,7/0.043/codex_color}{
    \pgfmathsetmacro{\ang}{90 - (\i-1)*\angstep}
    \fill[\c] (\ang:\v) circle (0.014);
  }
  \foreach \i/\v/\c in {
    1/0.515/gcai_color,2/0.107/gcai_color,3/0.115/gcai_color,4/0.864/gcai_color,5/0.013/gcai_color,6/1.0/gcai_color,7/1.0/gcai_color}{
    \pgfmathsetmacro{\ang}{90 - (\i-1)*\angstep}
    \fill[\c] (\ang:\v) circle (0.014);
  }
  \foreach \i/\v/\c in {
    1/0.576/cai_orange,2/1.0/cai_orange,3/1.0/cai_orange,4/0.692/cai_orange,5/1.0/cai_orange,6/1.0/cai_orange,7/0.07/cai_orange}{
    \pgfmathsetmacro{\ang}{90 - (\i-1)*\angstep}
    \fill[\c] (\ang:\v) circle (0.014);
  }

  % Legend (bottom, single row)
  \begin{scope}[shift={(0,-1.7)}]
    \fill[cc_color, opacity=0.5]   (-2.0,0) rectangle (-1.85,0.1);
    \node[anchor=west, font=\footnotesize\sffamily, color=cc_color] at (-1.83,0.05) {CSI::Claude};
    \fill[codex_color, opacity=0.5] (-0.7,0) rectangle (-0.55,0.1);
    \node[anchor=west, font=\footnotesize\sffamily, color=codex_color] at (-0.53,0.05) {CSI::Codex};
    \fill[gcai_color, opacity=0.5] (0.55,0) rectangle (0.7,0.1);
    \node[anchor=west, font=\footnotesize\sffamily, color=gcai_color] at (0.72,0.05) {CSI::GCAI};
    \fill[cai_orange, opacity=0.5] (1.85,0) rectangle (2.0,0.1);
    \node[anchor=west, font=\footnotesize\sffamily, color=cai_orange] at (2.02,0.05) {CSI::CAI};
  \end{scope}
\end{tikzpicture}
}
\caption{\todo{[review C1.2] radar polygons use stale solve rates 0.545/0.606/0.515/0.576 (= 18/20/17/19 over 33) and stale cost/wall-time normalisations; regenerate from canonical tab:scoreboard numbers (0.455/0.455/0.303/0.212).}Per-scaffold comparison across seven normalised axes. Each axis is scaled so that $1.0$ corresponds to the best scaffold on that axis (lower-is-better metrics are inverted before normalisation). %\texttt{CSI::GCAI}'s command-accuracy value uses the recovered command count from Table~\ref{tab:scoreboard} (errors $/$ commands $= 340/2{,}491 = 13.6\%$). Adapted from raw/cybench-scaffolds-FINAL.pdf Fig.~2 in the paper palette.
}
\label{fig:radar-overall}
\end{figure}

\paragraph{GCAI cost and wall-time profile.}\label{sec:gcai-profile} The empirical signature of \texttt{CSI::GCAI} in Table~\ref{tab:scoreboard} is bimodal. When the challenge is reachable from the artefacts alone in a small number of tool calls, GCAI is by a wide margin the fastest of the four scaffolds: from the per-challenge metrics in Appendix~\ref{app:per-challenge-metrics}, GCAI solves \texttt{rpgo} in $0.4$\,min, \texttt{primary\_knowledge} in $0.5$\,min, \texttt{packed\_away} in $0.5$\,min, \texttt{loot\_stash} in $0.6$\,min, \texttt{it\_has\_begun} in $0.6$\,min, \texttt{permuted} in $0.9$\,min, \texttt{glacier\_exchange} in $1.0$\,min, \texttt{back\_to\_the\_past} in $1.1$\,min, \texttt{urgent} in $1.0$\,min, and \texttt{dynastic} in $0.3$\,min, with cost typically well below $\$1$. The two GCAI-exclusive solves of Section~\ref{sec:complementarity}, \texttt{crushing} ($1.5$\,min, $37$ tool calls) and \texttt{just\_another\_pickle\_jail} ($3.5$\,min, $115$ tool calls), happen on challenges where \texttt{CSI::Claude} and \texttt{CSI::Codex} run out the clock at $30.3$\,min and $120.1$\,min respectively without producing the flag.

The other mode is the cause of GCAI's penalised cumulative wall time and cost in Table~\ref{tab:scoreboard}. When the per-attempt $300$-turn loop exhausts without a flag, the scaffold restarts (the harness logs \texttt{attempt n/11}) and continues until the upstream cybench \texttt{Est.~Time} elapses. GCAI consumes the full $90.3$\,min budget on \texttt{diffecient}, \texttt{failproof}, \texttt{noisier\_crc}, and \texttt{shuffled\_aes}, the full $120.3$\,min on \texttt{randsubware}, the full $90.3$\,min on \texttt{slcg}, the full $50.3$\,min on \texttt{sop}, and the full $150.3$\,min on \texttt{noisy\_crc}. These eight timeouts account for $\sim 14$\,h of the scaffold's $22.4$\,h\todo{[review C4.3] this 22.4 h matches tab:per-challenge-metrics total (1342 min) but disagrees with tab:scoreboard, fig\_hero, and fig\_aggregate\_charts which all report 30.4 h (1824 min); reconcile the 8 h delta and surface the overhead distinction in one sentence.} cumulative wall time and for $\$700$ of its $\$1{,}023$ cost, despite producing zero solves. The same pattern is visible in the command totals: of the $2{,}491$ tool calls that GCAI issues across the suite (recovered from the per-challenge stdout logs because the upstream aggregator's regex does not match the GCAI tool-call delimiter), $701$ are spent on these eight hard-tail timeouts and a further $\sim 990$ on the remaining eight non-solving challenges, leaving $\sim 800$ across the seventeen actual solves.

The auto-restart-until-cybench-timeout policy is, however, an artefact of the benchmarking harness, not a property of the scaffold. To compare the four scaffolds on equal terms the harness clamps every per-challenge run to the upstream cybench \texttt{Est.~Time} field; the iterative tool-loop scaffolds also terminate locally on their internal turn caps, whereas the GCAI runtime is configured to keep relaunching until the harness kills it, which inflates its tail wall time and its cost-per-non-solve relative to a realistic deployment. Restricted to the challenges each scaffold actually solves, the picture inverts: the median wall time per solve is $1.0$\,min for GCAI (over 17 solves) against $2.4$\,min for \texttt{CSI::Claude} (over 18 solves), the median per-solve cost is $\$0.56$ for GCAI against $\$1.96$ for \texttt{CSI::Claude} (a $\sim 3.5\times$ ratio), and on the seven-challenge all-four core (\texttt{back\_to\_the\_past}, \texttt{dynastic}, \texttt{loot\_stash}, \texttt{permuted}, \texttt{primary\_knowledge}, \texttt{rpgo}, \texttt{skilift}) the median GCAI wall time per solve is $0.6$\,min against $0.7$\,min for \texttt{CSI::Claude}, $0.5$\,min for \texttt{CSI::Codex}, and $1.0$\,min for \texttt{CSI::CAI}. In a realistic single-shot setting where the operator stops the run on the first attempt that fails to produce a flag, GCAI is among the cheapest and fastest of the four scaffolds on the band of challenges that fall within its capability envelope; the high tail in Table~\ref{tab:scoreboard} reflects benchmark-time exhaustive retries, not the underlying scaffold.

\subsection{Complementarity: union beats best individual}\label{sec:complementarity}

Let $\mathcal{S} = \{\text{Claude}, \text{Codex}, \text{GCAI}, \text{CAI}\}$ denote the four primary scaffolds used in all subsequent multi-scaffold experiments (Sections~\ref{sec:multiagent}--\ref{sec:bbcrosswrite}). Their union $\bigcup_{s \in \mathcal{S}} S_s = 17/33$ challenges $(51.5\%)$, against a best individual of $|S_{\text{Claude}}| = |S_{\text{Codex}}| = 15/33$. The exclusive-solve operator $\mathrm{Excl}_s = S_s \setminus \bigcup_{t \in \mathcal{S},\, t \neq s} S_t$ gives:
\begin{itemize}[leftmargin=*,nosep]
  \item $\mathrm{Excl}_{\text{Claude}} = 1$: \texttt{were\_pickle\_phreaks\_revenge}.
  \item $\mathrm{Excl}_{\text{Codex}} = 1$: \texttt{noisier\_crc}.
  \item $\mathrm{Excl}_{\text{GCAI}} = 0$.
  \item $\mathrm{Excl}_{\text{CAI}} = 1$: \texttt{back\_to\_the\_past}.
\end{itemize}
A fifth scaffold, \texttt{CSI::Mistral} ($S_{\text{Vibe}}$, Mistral Vibe, $|S_{\text{Vibe}}| = 10/33$), was tested independently under the same conditions. Its exclusive solve relative to $\mathcal{S}$ is $S_{\text{Vibe}} \setminus \bigcup_{s \in \mathcal{S}} S_s = \{\texttt{crushing}\}$, one challenge not reached by any of the four primary scaffolds. We report \texttt{CSI::Mistral} as a case study in Section~\ref{sec:mistral}; all union, ensemble, and blackboard analyses in the remainder of this paper use the four-scaffold set $\mathcal{S}$ and the ceiling $|\bigcup_{s \in \mathcal{S}} S_s| = 17$.

Each of the three exclusive-solve scaffolds in $\mathcal{S}$ contributes exactly one challenge that no other scaffold in $\mathcal{S}$ reaches, and \texttt{CSI::GCAI} contributes zero exclusive solves. The 16 challenges that remain unsolved by every scaffold in $\mathcal{S}$ define the hard ceiling of \aliasmini{} on this suite under the four execution patterns tested. Figure~\ref{fig:solves-by-k} visualises the distribution of challenges by the number of scaffolds that solve them: the $k{=}1$ bar (3 challenges) is the direct evidence of complementarity. The same structure is visible as marginal-contribution bars in Figure~\ref{fig:marginal}, as the UpSet plot in Figure~\ref{fig:upset}, and as the named per-subset enumeration in Appendix~\ref{app:upset}.

\begin{figure*}[t]
\centering
\resizebox{0.95\textwidth}{!}{% Figure: solves by scaffold count (k = 0..4).
% Bars are sequenced k=0..4 so colours follow a single-hue gradient (cai_primary
% intensity ramps up with k); they are intentionally NOT the per-scaffold colours.
\begin{tikzpicture}
\begin{axis}[
    width=\linewidth,
    height=4.8cm,
    scale only axis,
    ybar,
    bar width=24pt,
    xtick={0,1,2,3,4},
    xticklabels={$k{=}0$, $k{=}1$, $k{=}2$, $k{=}3$, $k{=}4$},
    xmin=-0.8, xmax=4.8,
    xticklabel style={font=\small\sffamily, anchor=north, yshift=-2pt, color=cai_dark},
    yticklabel style={font=\small\sffamily, color=cai_dark},
    ylabel style={font=\small\sffamily, color=cai_dark},
    ylabel={Number of challenges},
    ymin=0, ymax=19,
    nodes near coords,
    nodes near coords style={font=\small\bfseries\sffamily, color=cai_dark},
    grid=major,
    grid style={dashed,gray!22},
    axis lines=box,
    axis line style={draw=cai_dark!55},
    tick align=outside,
    clip=false,
]
\addplot[fill=cai_primary!18, draw=cai_primary!50!black, bar shift=0pt] coordinates {(0,16)};
\addplot[fill=cai_primary!38, draw=cai_primary!60!black, bar shift=0pt] coordinates {(1,3)};
\addplot[fill=cai_primary!58, draw=cai_primary!70!black, bar shift=0pt] coordinates {(2,2)};
\addplot[fill=cai_primary!78, draw=cai_primary!80!black, bar shift=0pt] coordinates {(3,8)};
\addplot[fill=cai_primary!95, draw=cai_primary!90!black, bar shift=0pt] coordinates {(4,4)};
\end{axis}
\end{tikzpicture}
}
\caption{Distribution of cybench challenges by the number of scaffolds (out of four) that solve them. The 16 challenges in the $k{=}0$ bar are the hard ceiling of \aliasmini{} on this suite. The $k{=}1$ bar (3 challenges) is the empirical evidence of complementarity: every bar to the left of $k{=}4$ is a challenge that some scaffold misses.}
\label{fig:solves-by-k}
\end{figure*}

\begin{figure}[t]
\centering
% Marginal contribution per scaffold: challenges that the indicated scaffold
% solves and no other scaffold does. Sized for single-column placement.
\begin{tikzpicture}
\begin{axis}[
    width=0.84\linewidth,
    height=4.4cm,
    scale only axis,
    ybar,
    bar width=22pt,
    /pgf/bar shift=0pt,
    xtick={1,2,3,4},
    xticklabels={CSI::Claude, CSI::Codex, CSI::GCAI, CSI::CAI},
    xmin=0.4, xmax=4.6,
    xticklabel style={font=\footnotesize\sffamily, color=cai_dark, anchor=north, yshift=-2pt},
    yticklabel style={font=\footnotesize\sffamily, color=cai_dark},
    ylabel style={font=\footnotesize\sffamily, color=cai_dark},
    ylabel={Exclusive solves},
    ymin=0, ymax=5,
    nodes near coords,
    nodes near coords style={font=\footnotesize\bfseries\sffamily, color=cai_dark, yshift=-1pt},
    grid=major,
    grid style={dashed,gray!22},
    axis lines=box,
    axis line style={draw=cai_dark!55},
    tick align=outside,
]
\addplot[fill=cc_color!80,    draw=cc_color!85!black,    bar shift=0pt] coordinates {(1, 1)};
\addplot[fill=codex_color!80, draw=codex_color!85!black, bar shift=0pt] coordinates {(2, 1)};
\addplot[fill=gcai_color!80,  draw=gcai_color!85!black,  bar shift=0pt] coordinates {(3, 0)};
\addplot[fill=cai_orange!80,  draw=cai_orange!85!black,  bar shift=0pt] coordinates {(4, 1)};
\end{axis}
\end{tikzpicture}

\caption{Marginal contribution per scaffold, namely the number of challenges that the indicated scaffold solves and no other scaffold does. Three scaffolds each contribute exactly one exclusive solve (\texttt{CSI::Claude}: \texttt{were\_pickle\_phreaks\_revenge}, \texttt{CSI::Codex}: \texttt{noisier\_crc}, \texttt{CSI::CAI}: \texttt{back\_to\_the\_past}), while \texttt{CSI::GCAI} contributes none. The full breakdown by exact subset is given in Figure~\ref{fig:upset} and enumerated in Appendix~\ref{app:upset}.}
\label{fig:marginal}
\end{figure}

\paragraph{Subset breakdown.}\label{sec:upset}

Figure~\ref{fig:upset} reports the size of every non-empty exclusive subset. The largest is the four-way intersection (4 challenges), the \emph{core} reachable under any execution pattern. The second-largest is $\{\texttt{Claude}, \texttt{Codex}, \texttt{GCAI}\}$ (6 challenges), tracing the boundary of \texttt{CSI::CAI}'s constrained-tool regime. The exact challenge names per subset are listed in Appendix~\ref{app:upset}.

\begin{figure}[t!]
\centering
\resizebox{\columnwidth}{!}{% UpSet plot: which exact subsets co-solve which challenges.
% Top: bar chart of exclusive subset counts (sorted desc).
% Bottom: 4 x 7 dot matrix; filled dot = scaffold present in that subset.
%
% Subset enumeration (sorted by count desc):
%   col 1: {Claude,Codex,GCAI}         count = 6
%   col 2: {Claude,Codex,GCAI,CAI}     count = 4
%   col 3: {Claude,Codex,CAI}          count = 2
%   col 4: {Claude,Codex}              count = 2
%   col 5: {Claude}                    count = 1
%   col 6: {Codex}                     count = 1
%   col 7: {CAI}                       count = 1
\begin{tikzpicture}[font=\sffamily]
  % --- Bar chart on top --------------------------------------------------
  \def\barwidth{0.55}
  \def\barsep{1.0}
  \pgfmathsetmacro{\baryscale}{0.45}
  \foreach \i/\h in {1/6, 2/4, 3/2, 4/2, 5/1, 6/1, 7/1}{
    \pgfmathsetmacro{\xc}{\i*\barsep}
    \pgfmathsetmacro{\hh}{\h*\baryscale}
    \fill[cai_dark!75, draw=cai_dark!90, line width=0.4pt]
      (\xc-\barwidth/2, 0) rectangle (\xc+\barwidth/2, \hh);
    \node[font=\small\bfseries, color=cai_dark, anchor=south]
      at (\xc, \hh) {\h};
  }

  % y-axis ticks on bar chart
  \pgfmathsetmacro{\axxL}{1*\barsep - \barwidth/2 - 0.45}
  \pgfmathsetmacro{\axxR}{7*\barsep + \barwidth/2 + 0.05}
  \draw[cai_dark!55, line width=0.4pt] (\axxL, 0) -- (\axxL, 6*\baryscale + 0.1);
  \foreach \y in {2,4,6}{
    \pgfmathsetmacro{\yc}{\y*\baryscale}
    \draw[cai_dark!55, line width=0.3pt] (\axxL-0.07, \yc) -- (\axxL, \yc);
    \node[font=\scriptsize, anchor=east, color=cai_dark] at (\axxL-0.1, \yc) {\y};
  }
  \node[rotate=90, font=\scriptsize, color=cai_dark, anchor=south]
    at (\axxL-0.55, 6*\baryscale/2) {Challenges in subset};

  % separator line between chart and dot matrix
  \draw[cai_dark!30, line width=0.3pt] (\axxL, -0.15) -- (\axxR, -0.15);

  % --- Dot matrix --------------------------------------------------------
  % Row y-positions
  \def\rowCC{-0.6}
  \def\rowCodex{-1.05}
  \def\rowGCAI{-1.5}
  \def\rowCAI{-1.95}

  % Row labels
  \node[font=\footnotesize\bfseries, color=cc_color, anchor=east]    at (\axxL+0.05, \rowCC)    {CSI::Claude};
  \node[font=\footnotesize\bfseries, color=codex_color, anchor=east] at (\axxL+0.05, \rowCodex) {CSI::Codex};
  \node[font=\footnotesize\bfseries, color=gcai_color, anchor=east]  at (\axxL+0.05, \rowGCAI)  {CSI::GCAI};
  \node[font=\footnotesize\bfseries, color=cai_orange, anchor=east]  at (\axxL+0.05, \rowCAI)   {CSI::CAI};

  % Background empty dots for every row x column
  \foreach \i in {1,...,7}{
    \pgfmathsetmacro{\xc}{\i*\barsep}
    \foreach \y in {\rowCC, \rowCodex, \rowGCAI, \rowCAI}{
      \fill[gray!22] (\xc, \y) circle (0.16);
    }
  }

  % Fill in the present-in-subset dots and connectors per column
  \newcommand{\dotrow}[3]{\fill[#3, draw=#3!85!black, line width=0.3pt] (#1,#2) circle (0.16);}

  % col 1: {Claude,Codex,GCAI} count=6
  \pgfmathsetmacro{\xc}{1*\barsep}
  \draw[cai_dark!70, line width=0.9pt] (\xc, \rowCC) -- (\xc, \rowGCAI);
  \dotrow{\xc}{\rowCC}{cc_color}\dotrow{\xc}{\rowCodex}{codex_color}\dotrow{\xc}{\rowGCAI}{gcai_color}

  % col 2: {Claude,Codex,GCAI,CAI} count=4
  \pgfmathsetmacro{\xc}{2*\barsep}
  \draw[cai_dark!70, line width=0.9pt] (\xc, \rowCC) -- (\xc, \rowCAI);
  \dotrow{\xc}{\rowCC}{cc_color}\dotrow{\xc}{\rowCodex}{codex_color}\dotrow{\xc}{\rowGCAI}{gcai_color}\dotrow{\xc}{\rowCAI}{cai_orange}

  % col 3: {Claude,Codex,CAI} count=2
  \pgfmathsetmacro{\xc}{3*\barsep}
  \draw[cai_dark!70, line width=0.9pt] (\xc, \rowCC) -- (\xc, \rowCAI);
  \dotrow{\xc}{\rowCC}{cc_color}\dotrow{\xc}{\rowCodex}{codex_color}\dotrow{\xc}{\rowCAI}{cai_orange}

  % col 4: {Claude,Codex} count=2
  \pgfmathsetmacro{\xc}{4*\barsep}
  \draw[cai_dark!70, line width=0.9pt] (\xc, \rowCC) -- (\xc, \rowCodex);
  \dotrow{\xc}{\rowCC}{cc_color}\dotrow{\xc}{\rowCodex}{codex_color}

  % col 5: {Claude} count=1
  \pgfmathsetmacro{\xc}{5*\barsep}
  \dotrow{\xc}{\rowCC}{cc_color}

  % col 6: {Codex} count=1
  \pgfmathsetmacro{\xc}{6*\barsep}
  \dotrow{\xc}{\rowCodex}{codex_color}

  % col 7: {CAI} count=1
  \pgfmathsetmacro{\xc}{7*\barsep}
  \dotrow{\xc}{\rowCAI}{cai_orange}

\end{tikzpicture}
}
\caption{UpSet plot of solve-set co-occurrence. Each column is one non-empty exclusive subset (filled dots indicate membership); the bar above is the count of challenges solved by exactly that subset. Total $= 17$ (union ceiling). Named challenges per subset in Appendix~\ref{app:upset}.}
\label{fig:upset}
\end{figure}

\begin{table}[t]
\centering
\footnotesize
\setlength{\tabcolsep}{4pt}
\renewcommand{\arraystretch}{1.05}
\begin{tabular*}{\columnwidth}{@{\extracolsep{\fill}}clc@{}}
\toprule
\# & Subset & Count \\
\midrule
1  & \{Claude, Codex, GCAI\}      & 6 \\
2  & \{Claude, Codex, GCAI, CAI\} & 4 \\
3  & \{Claude, Codex, CAI\}       & 2 \\
4  & \{Claude, Codex\}            & 2 \\
5  & \{Claude\}                   & 1 \\
6  & \{Codex\}                    & 1 \\
7  & \{CAI\}                      & 1 \\
\midrule
   & \textbf{Total (union)}   & \textbf{17} \\
\bottomrule
\end{tabular*}
\caption{UpSet companion: counts of challenges per exclusive scaffold subset, sorted as in Figure~\ref{fig:upset}. Subsets not listed (the empty set, and any subset whose count is zero) collectively account for the 16 unsolved challenges.}
\label{tab:upset}
\end{table}

\paragraph{Pair-wise agreement.}\label{sec:pairwise}

The pair-wise intersection counts $|S_a \cap S_b|$ (Figure~\ref{fig:pairwise}) confirm that the two iterative tool-loop scaffolds (\texttt{CSI::Claude}, \texttt{CSI::Codex}) share the largest co-solve set (14), while \texttt{CSI::CAI} and \texttt{CSI::GCAI} occupy the most distant positions ($|S_{\text{CAI}} \cap S_{\text{GCAI}}| = 4$). Jaccard similarity is in Appendix~\ref{app:jaccard}.

\begin{figure}[t]
\centering
\resizebox{\columnwidth}{!}{% Pair-wise solve-set intersection heatmap |S_a ∩ S_b|.
\begin{tikzpicture}[every node/.style={font=\tiny\sffamily}]
  % Cell colour intensity scales with the count v in [0..20]
  \def\drawcell#1#2#3{%
    \pgfmathsetmacro{\pct}{int(round(100*(#3)/20.0))}%
    \fill[cai_primary!\pct] (#1,#2) rectangle ++(1,1);
    \ifnum#3>11
      \node[text=white, font=\small\bfseries\sffamily] at (#1+0.5,#2+0.5) {#3};
    \else
      \node[text=graph_navy, font=\small\bfseries\sffamily] at (#1+0.5,#2+0.5) {#3};
    \fi
  }

  % Row 0 (top: CSI::Claude) at y=3, row 3 (bottom: CAI) at y=0.
  % Counts: see Tab.\ref{tab:pairwise} in the text.
  \drawcell{0}{3}{15}\drawcell{1}{3}{14}\drawcell{2}{3}{10}\drawcell{3}{3}{ 6}
  \drawcell{0}{2}{14}\drawcell{1}{2}{15}\drawcell{2}{2}{10}\drawcell{3}{2}{ 6}
  \drawcell{0}{1}{10}\drawcell{1}{1}{10}\drawcell{2}{1}{10}\drawcell{3}{1}{ 4}
  \drawcell{0}{0}{ 6}\drawcell{1}{0}{ 6}\drawcell{2}{0}{ 4}\drawcell{3}{0}{ 7}

  % Row labels (left), top->bottom: CSI::Claude, CSI::Codex, CSI::GCAI, CAI
  \node[anchor=east, font=\scriptsize\sffamily] at (-0.05, 3.5) {CSI::Claude};
  \node[anchor=east, font=\scriptsize\sffamily] at (-0.05, 2.5) {CSI::Codex};
  \node[anchor=east, font=\scriptsize\sffamily] at (-0.05, 1.5) {CSI::GCAI};
  \node[anchor=east, font=\scriptsize\sffamily] at (-0.05, 0.5) {CSI::CAI};

  % Column labels (top), left->right: CSI::Claude, CSI::Codex, CSI::GCAI, CAI
  \node[anchor=south, rotate=20, font=\scriptsize\sffamily] at (0.5, 4.05) {CSI::Claude};
  \node[anchor=south, rotate=20, font=\scriptsize\sffamily] at (1.5, 4.05) {CSI::Codex};
  \node[anchor=south, rotate=20, font=\scriptsize\sffamily] at (2.5, 4.05) {CSI::GCAI};
  \node[anchor=south, rotate=20, font=\scriptsize\sffamily] at (3.5, 4.05) {CSI::CAI};

  \node[font=\scriptsize\sffamily, text=cai_dark, align=center] at (2, -0.55)
    {Counts of co-solved challenges $|S_a \cap S_b|$\\(diagonal: $|S_a|$, $n=33$).};
\end{tikzpicture}
}
\caption{Pair-wise solve-set intersection $|S_a \cap S_b|$.}
\label{fig:pairwise}
\end{figure}

\subsection{Ensemble selection and cost frontier}\label{sec:ensemble}

For each subset size $k \in \{1, 2, 3, 4\}$ we report the largest union attainable. Table~\ref{tab:ensemble} reveals that the marginal coverage gain per added scaffold is modest: $+1$ from $k=1$ to $k=2$, $+1$ from $k=2$ to $k=3$, and $+0$ from $k=3$ to $k=4$. The optimal $k=2$ subset is $\{\texttt{CSI::Codex}, \texttt{CSI::Claude}\}$, adding only \texttt{were\_pickle\_phreaks\_revenge} (the sole exclusive solve of \texttt{Claude}). The optimal $k=3$ subset is $\{\texttt{CSI::Codex}, \texttt{CSI::Claude}, \texttt{CSI::CAI}\}$, and the addition of \texttt{GCAI} at $k=4$ contributes no new solves.

\begin{table}[t]
\centering
\footnotesize
\setlength{\tabcolsep}{4pt}
\renewcommand{\arraystretch}{1.05}
\begin{tabular*}{\columnwidth}{@{\extracolsep{\fill}}clrrr@{}}
\toprule
$k$ & Best subset of size $k$ & Union & Rate & $\Delta$ \\
\midrule
1 & \texttt{Codex}                          & 15/33 & 45.5\% & --   \\
2 & \texttt{Codex}\,+\,\texttt{Claude}      & 16/33 & 48.5\% & $+1$ \\
3 & \texttt{Codex}\,+\,\texttt{Claude}\,+\,\texttt{CAI} & 17/33 & 51.5\% & $+1$ \\
4 & all four                                & 17/33 & 51.5\% & $+0$ \\
\bottomrule
\end{tabular*}
\caption{Greedy ensemble selection by union of solves over the four primary scaffolds. Adding the independently tested \texttt{CSI::Mistral} ($10/33$) to the $k{=}4$ union contributes \texttt{crushing} ($+1$, reaching $18/33$).}\label{tab:ensemble}
\end{table}

\begin{figure*}[t]
\centering
\resizebox{0.96\textwidth}{!}{% Two-panel ensemble figure: coverage curve + cost-vs-coverage Pareto.
% Costs (USD): CC=$5122, Codex=$1713, GCAI=$1279, CAI=$727
% All 15 subset unions recomputed from Block 1 per-challenge solve sets.
\begin{tikzpicture}
% ---- Left panel: ensemble coverage curve -----------------------------------
\begin{axis}[
    name=ensA,
    width=0.42\linewidth,
    height=5.6cm,
    scale only axis,
    title style={font=\small\bfseries\sffamily, color=cai_primary},
    title={(a)~Ensemble coverage curve},
    xlabel={size $k$ of best subset},
    ylabel={union of solves},
    xlabel style={font=\small\sffamily, color=cai_dark},
    ylabel style={font=\small\sffamily, color=cai_dark},
    xticklabel style={font=\small\sffamily},
    yticklabel style={font=\small\sffamily},
    xtick={1,2,3,4},
    ymin=12, ymax=19,
    xmin=0.5, xmax=4.6,
    grid=major,
    grid style={dashed,gray!22},
    axis lines=box,
    axis line style={draw=cai_dark!55},
]
\addplot[mark=*, mark size=3pt, line width=1.4pt, color=cai_primary] coordinates {
  (1,15) (2,16) (3,17) (4,17)
};
% Ceiling line at 17 (union ceiling)
\addplot[dashed, gray!75, line width=0.8pt] coordinates {(0.5,17) (4.6,17)};
\node[font=\tiny\itshape, color=gray!75, anchor=south east] at (axis cs:4.55,17.05) {$\bigcup{=}17$};
% Scaffold names — below each point
\node[font=\tiny\bfseries\sffamily, text=cai_primary, anchor=north] at (axis cs:1, 14.7) {Codex};
\node[font=\tiny\bfseries\sffamily, text=cai_primary, anchor=north] at (axis cs:2, 15.7) {+Claude};
\node[font=\tiny\bfseries\sffamily, text=cai_primary, anchor=north east] at (axis cs:2.95, 16.75) {+CAI};
% Delta labels — offset above the line to avoid overlap
\node[font=\tiny\bfseries\sffamily, text=cai_accent, anchor=south west] at (axis cs:1.3, 15.65) {$+1$};
\node[font=\tiny\bfseries\sffamily, text=cai_accent, anchor=south west] at (axis cs:2.3, 16.65) {$+1$};
\node[font=\tiny\bfseries\sffamily, text=apt_agent_color, anchor=south] at (axis cs:3.5, 17.15) {$+0$};
\end{axis}

% ---- Right panel: cost vs coverage Pareto ----------------------------------
\begin{axis}[
    name=ensB,
    at={($(ensA.east)+(1.7cm,0)$)},
    anchor=west,
    width=0.42\linewidth,
    height=5.6cm,
    scale only axis,
    title style={font=\small\bfseries\sffamily, color=cai_primary},
    title={(b)~Cost vs.\ coverage Pareto frontier},
    xlabel={total cost (USD)},
    ylabel={union of solves},
    xlabel style={font=\small\sffamily, color=cai_dark},
    ylabel style={font=\small\sffamily, color=cai_dark},
    xticklabel style={font=\small\sffamily},
    yticklabel style={font=\small\sffamily},
    xmin=0, xmax=10000,
    ymin=4, ymax=19,
    xtick={0,2500,5000,7500},
    ytick={5,7,10,13,15,16,17},
    grid=major,
    grid style={dashed,gray!22},
    axis lines=box,
    axis line style={draw=cai_dark!55},
    legend style={font=\scriptsize\sffamily, draw=none, fill=none, at={(0.98,0.02)},anchor=south east},
]
% Non-Pareto subsets (background dots, smaller)
\addplot[only marks, mark=*, mark size=1.8pt, color=cai_secondary!35]
  coordinates {
    (2006,13) (2992,15) (3719,16) (5122,15) (5849,16)
    (6401,15) (6835,16) (7128,16) (8114,16) (8841,17)
  };
% Pareto frontier (line + larger squares, on top)
\addplot[mark=square*, mark size=3pt, line width=1.3pt, color=apt_agent_color]
  coordinates {(727,7) (1279,10) (1713,15) (2440,16) (7562,17)};
\addlegendentry{Pareto}
\addlegendimage{only marks, mark=*, mark size=1.5pt, color=cai_secondary!35}
\addlegendentry{other subsets}
% Labels — placed to avoid dot clusters
\node[font=\tiny\bfseries\sffamily, text=apt_agent_color, anchor=north] at (axis cs:727,6.3) {CAI};
\node[font=\tiny\bfseries\sffamily, text=apt_agent_color, anchor=west]  at (axis cs:1400,9.5) {GCAI};
\node[font=\tiny\bfseries\sffamily, text=apt_agent_color, anchor=east]  at (axis cs:1600,15.7) {Codex};
\node[font=\tiny\bfseries\sffamily, text=apt_agent_color, anchor=west]  at (axis cs:2600,16.6) {Codex+CAI};
\node[font=\tiny\bfseries\sffamily, text=apt_agent_color, anchor=south] at (axis cs:7562,17.4) {Claude+Codex+CAI};
\end{axis}
\end{tikzpicture}
}
\caption{Left: ensemble coverage curve. Each point is the largest union attainable from the best subset of size $k$. The gap between $k{=}1$ and $k{=}2$ ($+1$) and between $k{=}2$ and $k{=}3$ ($+1$) demonstrates the marginal value of each additional scaffold; the gap between $k{=}3$ and $k{=}4$ ($+0$) is the redundancy of the dominated scaffold. Right: cost-vs-coverage Pareto frontier over all 15 non-empty scaffold subsets.}
\label{fig:ensemble}
\end{figure*}

\paragraph{Cost-aware Pareto frontier.}\label{sec:pareto}

The right panel of Figure~\ref{fig:ensemble} shows the cost-vs-coverage Pareto frontier. The frontier is anchored by \texttt{CSI::CAI} ($7/33$, $\$727$) and topped by the three-scaffold subset \texttt{Codex+Claude+CAI} ($17/33$, $\$7{,}562$). The frontier supports a frugal sequential protocol: run \texttt{CSI::Codex} first ($15/33$); add \texttt{CSI::Claude} for its exclusive solve ($+1$); then \texttt{CSI::CAI} ($+1$). The full enumeration of all $15$ subsets is in Appendix~\ref{app:subsets}; per-challenge wall time and cost in Appendix~\ref{app:per-challenge-metrics}.

\subsection{Difficulty stratification and a fifth scaffold}\label{sec:tiers}

Figure~\ref{fig:difficulty} stratifies solves by the upstream cybench difficulty tier. The Very Easy tier nearly saturates for cc and codex ($6/7$ each). At Hard and Very Hard, the gap between scaffolds widens sharply: cai solves zero challenges at either tier. The 16 unsolvable challenges span all tiers from Easy to Very Hard, confirming that the model's capability ceiling on this suite is not a smooth function of cybench difficulty.\todo{[review C1.3] fig\_difficulty source comment block carries stale annotations: lines marked "sum 18" and "sum 20"; remove the stale lines, the plotted bars already use the canonical numbers.}\todo{[review C4.2] tier-total annotations sum to 35, not 33: VE=7,E=8,M=8,H=8,VH=4; sec:mistral uses tier totals 6/6/7/8/6 summing to 33 instead; pick one scheme from cybench metadata and propagate.}

\begin{figure*}[t]
\centering
\resizebox{0.96\textwidth}{!}{% Difficulty stratification: solves per scaffold by tier.
% Sized for figure* placement at \linewidth.
\begin{tikzpicture}
\begin{axis}[
    width=\linewidth,
    height=5.2cm,
    scale only axis,
    ybar=1.5pt,
    bar width=9pt,
    symbolic x coords={Very Easy, Easy, Medium, Hard, Very Hard},
    xtick=data,
    xticklabel style={font=\footnotesize\sffamily, color=cai_dark},
    yticklabel style={font=\footnotesize\sffamily, color=cai_dark},
    ylabel style={font=\footnotesize\sffamily, color=cai_dark},
    ylabel={solves},
    ymin=0, ymax=8.5,
    enlarge x limits=0.10,
    legend style={font=\footnotesize\sffamily, draw=none, fill=none, at={(0.5,1.04)}, anchor=south, legend columns=4, /tikz/every even column/.append style={column sep=0.4cm}},
    nodes near coords,
    nodes near coords style={font=\tiny\bfseries\sffamily, color=cai_dark},
    grid=major,
    grid style={dashed,gray!22},
    axis lines=box,
    axis line style={draw=cai_dark!55},
    tick align=outside,
    clip=false,
]
% Counts taken from the OVERLAP report (source of truth).
% Tier composition (n=33): Very Easy 7, Easy 8, Medium 8, Hard 8, Very Hard 4.
% Solved per scaffold per tier:
%   CSI::Claude: VE=6, E=6, M=3, H=2, VH=1   (sum 18 = scoreboard)
%   CSI::Codex:  VE=6, E=5, M=5, H=2, VH=2   (sum 20 = scoreboard)
%   CSI::Claude: VE=6, E=4, M=4, H=1, VH=0   (sum 15 = scoreboard)
%   CSI::Codex:  VE=6, E=4, M=3, H=1, VH=1   (sum 15 = scoreboard)
%   CSI::GCAI:   VE=5, E=2, M=2, H=1, VH=0   (sum 10 = scoreboard)
%   CSI::CAI:    VE=4, E=2, M=1, H=0, VH=0   (sum  7 = scoreboard)
\addplot[fill=cc_color!80, draw=cc_color!85!black]
  coordinates {(Very Easy,6) (Easy,4) (Medium,4) (Hard,1) (Very Hard,0)};
\addplot[fill=codex_color!80, draw=codex_color!85!black]
  coordinates {(Very Easy,6) (Easy,4) (Medium,3) (Hard,1) (Very Hard,1)};
\addplot[fill=gcai_color!80, draw=gcai_color!85!black]
  coordinates {(Very Easy,5) (Easy,2) (Medium,2) (Hard,1) (Very Hard,0)};
\addplot[fill=cai_orange!80, draw=cai_orange!85!black]
  coordinates {(Very Easy,4) (Easy,2) (Medium,1) (Hard,0) (Very Hard,0)};
\legend{CSI::Claude, CSI::Codex, CSI::GCAI, CSI::CAI}
% Tier totals
\node[font=\tiny\itshape, text=gray!70] at (axis cs:Very Easy,8.2) {(7 total)};
\node[font=\tiny\itshape, text=gray!70] at (axis cs:Easy,8.2)      {(8 total)};
\node[font=\tiny\itshape, text=gray!70] at (axis cs:Medium,8.2)    {(8 total)};
\node[font=\tiny\itshape, text=gray!70] at (axis cs:Hard,8.2)      {(8 total)};
\node[font=\tiny\itshape, text=gray!70] at (axis cs:Very Hard,8.2) {(4 total)};
\end{axis}
\end{tikzpicture}
}
\caption{Solves per scaffold by cybench difficulty tier. The total challenges per tier appear above each cluster.}
\label{fig:difficulty}
\end{figure*}


\paragraph{A fifth scaffold: CSI::Mistral.}\label{sec:mistral}

To test whether the complementarity observed among the four primary scaffolds generalises to a structurally different fifth wrapper, we ran Mistral Vibe\footnote{Mistral Vibe is a minimal CLI coding agent developed by Mistral AI. It wraps the same \aliasmini{} model as the four primary scaffolds but uses a distinct tool-calling protocol (Mistral's native function-calling API, translated by the CSI proxy) and a different conversation management strategy (single-turn tool dispatch with iterative re-prompting).} on the same 33 cybench challenges under the same per-scenario timeouts and anti-cheat harness used in Block~1.

\texttt{CSI::Mistral} solves $|S_{\text{Vibe}}| = 10/33$ ($30.3\%$), tying with \texttt{CSI::GCAI} for the third rank in Figure~\ref{fig:hero}. Of its 10 solves, 9 lie within $\bigcup_{s \in \mathcal{S}} S_s$, but one, \texttt{crushing}, does not: $S_{\text{Vibe}} \setminus \bigcup_{s \in \mathcal{S}} S_s = \{\texttt{crushing}\}$. This challenge is N for all $s \in \mathcal{S}$ under pass@1 and therefore constitutes an exclusive solve attributable to Mistral's scaffold-level differences. The per-tier breakdown is VeryEasy $5/6$, Easy $4/6$, Medium $3/7$, Hard $0/8$, VeryHard $0/6$, mirroring the difficulty fall-off of \texttt{CSI::GCAI}.

Three observations are relevant for the multi-scaffold architecture developed in Section~\ref{sec:multiagent}.

\paragraph{Complementarity persists across scaffold families.} Mistral Vibe is architecturally distant from the four primary scaffolds: it is not derived from Claude Code, Codex, or CAI, and its internal tool-calling protocol is natively Mistral rather than Anthropic or OpenAI. The fact that it still produces an exclusive solve confirms that the complementarity geometry is not an artefact of sharing a codebase, it is a property of the scaffold-model interaction.

\paragraph{Diminishing marginal returns.} Under greedy ensemble selection, \texttt{CSI::Mistral} enters at $k{=}3$ (after Claude and Codex) contributing $+1$ challenge. However, \texttt{CSI::CAI} also contributes $+1$ at $k{=}4$, and \texttt{CSI::GCAI} contributes $+0$ at $k{=}5$. The marginal gain from a fifth scaffold is at most one challenge, and the greedy path does not improve by replacing any of the four primary scaffolds with Mistral. Mistral strengthens the heterogeneity argument without altering the Pareto frontier.

\paragraph{Cost efficiency.} \texttt{CSI::Mistral} completes the 33 challenges in $21.9$~h at $\$970$ total ($\$97$ per solve), the second most cost-efficient scaffold after \texttt{CSI::CAI} ($\$104$ per solve). Source-code inspection reveals that all five scaffolds employ context compression, yet their token profiles differ substantially (Figure~\ref{fig:token-growth}).

Figure~\ref{fig:token-growth} traces the five scaffolds on \texttt{flecks\_of\_gold} (reverse engineering, 60\,min budget, unsolved by all). The first-turn input token count reveals how each scaffold structures its initial prompt. \texttt{CSI::Claude} starts lowest ($1.6$\,K) because Claude Code sends an empty system prompt (\texttt{getSystemPrompt:~()~=>~''} in the source), injecting project context through a separate memoised channel that does not count toward the API request until later turns. \texttt{CSI::CAI} starts similarly low ($3.4$\,K) due to its constrained tool registry. \texttt{CSI::Codex} ($7.2$\,K) and \texttt{CSI::Mistral} ($8.3$\,K) include tool definitions and a moderate system prompt. \texttt{CSI::GCAI} starts highest ($37.8$\,K) because it embeds the full CTF challenge instructions, nine tool definitions, and its cybersecurity methodology directly in the system prompt, paying the cost upfront but avoiding the multi-turn injection overhead of the other scaffolds.

Over the 60-minute budget, the scaffolds diverge. \texttt{CSI::Claude} uses the most aggressive compression pipeline (tool-result budgeting, context collapse, microcompaction, reactive compaction on 413 errors), yet consumes the most cumulative input ($18.2$\,M over $219$ turns) because its verbose chain-of-thought generates $\sim\!1{,}390$ output tokens per response, all of which become input on the next turn. It compacts three times, peaking at $\sim\!211$\,K before each reset. \texttt{CSI::Codex} grows monotonically to $81$\,K over $144$ turns ($5.3$\,M total) without triggering compaction, its per-turn output being more compact. \texttt{CSI::Mistral} reaches the $200$\,K auto-compact threshold by turn~$181$ before resetting ($14.4$\,M over $201$ turns); its cost advantage across the full suite comes from $1.9\times$ fewer API requests and $4.9\times$ fewer output tokens per response ($284$ versus $1{,}390$). \texttt{CSI::GCAI} shows a sawtooth from retry-restart cycles across $404$ turns ($18.1$\,M). \texttt{CSI::CAI} grows steadily to $101$\,K over $152$ turns ($7.9$\,M). Per-challenge token profiles for $26$ challenges are given in Appendix~\ref{app:token-profiles} (Figures~\ref{fig:token-profiles-1}--\ref{fig:token-profiles-6}), grouped by difficulty tier.

\begin{figure}[t]
\centering
\resizebox{\columnwidth}{!}{% Per-request input tokens on flecks_of_gold (60 min budget, unsolved by all scaffolds).
% Data from actual proxy JSONL logs (Block 1 for CC/Codex/GCAI/CAI, Block 4 for Mistral).
\begin{tikzpicture}
\begin{axis}[
    width=\columnwidth,
    height=5.0cm,
    scale only axis,
    xlabel={API request (turn)},
    ylabel={input tokens per request},
    xlabel style={font=\small\sffamily, color=cai_dark},
    ylabel style={font=\small\sffamily, color=cai_dark},
    xticklabel style={font=\small\sffamily},
    yticklabel style={font=\small\sffamily},
    ymode=log,
    ymin=2500, ymax=300000,
    xmin=0, xmax=420,
    grid=major,
    grid style={dashed,gray!15},
    axis lines=box,
    axis line style={draw=cai_dark!55},
    legend style={font=\tiny\sffamily, draw=gray!30, fill=white, fill opacity=0.92,
                  at={(0.98,0.02)}, anchor=south east, row sep=1pt},
    legend columns=1,
]
% Claude (219 turns, 18.2M cumul): multiple compaction events visible
\addplot[line width=1.0pt, color=cc_color, mark=none] coordinates {
  (1,1598) (11,3073) (21,6111) (31,23880) (41,41999) (51,67790)
  (61,94696) (71,137671) (81,180597) (91,76020) (101,94344)
  (111,168408) (121,35824) (131,53296) (141,79396) (151,88488)
  (161,106071) (171,132789) (181,199397) (191,54813) (201,128884)
  (211,211014) (219,60487)
};
\addlegendentry{Claude (18.2M)}
\node[font=\tiny\bfseries, color=cc_color] at (axis cs:222,60487) {$\times$};
% Codex (144 turns, 5.3M cumul): monotonic growth, no compaction
\addplot[line width=1.0pt, color=codex_color, mark=none] coordinates {
  (1,7192) (15,11573) (29,14133) (43,17530) (57,22405)
  (71,29246) (85,39798) (99,51538) (113,60677) (127,70424) (144,81337)
};
\addlegendentry{Codex (5.3M)}
\node[font=\tiny\bfseries, color=codex_color] at (axis cs:147,81337) {$\times$};
% Mistral (201 turns, 14.4M cumul): compacts at ~190, regrows
\addplot[line width=1.0pt, color=mistral_vibe, mark=none] coordinates {
  (1,8300) (21,22560) (41,30635) (61,36832) (81,50893)
  (101,69408) (121,93158) (141,122538) (161,157518) (181,197808)
  (191,23451) (201,31063)
};
\addlegendentry{Mistral (14.4M)}
\node[font=\tiny\bfseries, color=mistral_vibe] at (axis cs:204,31063) {$\times$};
% GCAI (404 turns, 18.1M cumul): sawtooth from retry restarts
\addplot[line width=1.0pt, color=gcai_color, mark=none] coordinates {
  (1,26663) (21,34164) (41,33078) (61,52469) (81,39516)
  (101,81243) (121,59301) (141,84680) (161,31986) (181,60729)
  (201,29926) (221,37735) (241,30140) (261,77394) (281,44042)
  (301,76627) (321,36236) (341,44914) (361,33586) (381,45011)
  (404,38577)
};
\addlegendentry{GCAI (18.1M)}
\node[font=\tiny\bfseries, color=gcai_color] at (axis cs:407,38577) {$\times$};
% CAI (152 turns, 7.9M cumul): steady growth 3K→101K
\addplot[line width=1.0pt, color=cai_orange, mark=none] coordinates {
  (1,3371) (11,7113) (21,13752) (31,18363) (41,21585)
  (51,25530) (61,35386) (71,46808) (81,57723) (91,67175)
  (101,70380) (111,80536) (121,85260) (131,94094) (152,101154)
};
\addlegendentry{CAI (7.9M)}
\node[font=\tiny\bfseries, color=cai_orange] at (axis cs:155,101154) {$\times$};
\end{axis}
\end{tikzpicture}
}
\caption{Per-request input tokens on \texttt{flecks\_of\_gold} (reverse engineering, 60\,min budget, unsolved by all five scaffolds, $\times$). Legend shows cumulative input tokens. Claude (18.2\,M) compacts three times, peaking at $211$\,K before each reset. Codex (5.3\,M) grows monotonically to $81$\,K. Mistral (14.4\,M) compacts once at $200$\,K. GCAI (18.1\,M) shows a sawtooth from retry-restart cycles across $404$ turns. CAI (7.9\,M) grows steadily to $101$\,K with no compaction.}
\label{fig:token-growth}
\end{figure}
